
\subsection{Random}

The random agent is the simplest baseline in our experiments. At each decision
point, the environment provides the agent with an information set \(I\),
including the current player's hand, the last move, public history features, and
the legal action set \(\mathcal{A}(I)\). The random agent ignores all strategic
information and samples uniformly from the legal actions:

\[
    a \sim \mathrm{Uniform}(\mathcal{A}(I)).
\]

In implementation, this corresponds directly to
\texttt{random.choice(infoset.legal\_actions)}. Since the action is always
sampled from the environment-generated legal action list, the agent never needs
to check card-type validity by itself. This baseline is useful because it gives
a lower-bound reference for whether a learned or heuristic policy is making
meaningful decisions beyond simply playing legal cards.